\chapter{Explicit Effects}\label{chap:explicit-effects}

% Notation of this chapter.
\newcommand{\Tok}[1]{\mathsf{heap}_{#1}}
\newcommand{\Rest}{\mathsf{rest}}
\newcommand{\getf}[2]{\mathsf{get}^{#2}_{#1}}
\newcommand{\setf}[2]{\mathsf{set}^{#2}_{#1}}
\newcommand{\initf}[2]{\mathsf{init}^{#2}_{#1}}
\newcommand{\newc}[1]{\mathsf{new}_{#1}}
\newcommand{\guard}{\mathsf{guard}}
\newcommand{\dup}{\mathsf{dup}}
\newcommand{\drop}{\mathsf{drop}}
\newcommand{\move}{\mathsf{move}}
\newcommand{\jump}{\mathsf{jump}}
\newcommand{\kif}[3]{\mathsf{if}\ #1\ \mathsf{then}\ #2\ \mathsf{else}\ #3}
\newcommand{\kret}{\mathsf{return}}
\newcommand{\ktrap}{\mathsf{trap}}
\newcommand{\Ret}{\mathsf{ret}}
\newcommand{\Exc}{\mathsf{exc}}
\newcommand{\Trap}{\lightning}
\newcommand{\Loc}{\mathit{Loc}}
\newcommand{\Obj}{\mathit{Obj}}
\newcommand{\Val}{\mathit{Val}}
\newcommand{\AObj}{\mathit{AObj}}
\newcommand{\alloc}{\mathsf{alloc}}
\newcommand{\pts}{\mathit{pts}}
\newcommand{\lab}{\mathit{label}}
\newcommand{\callees}{\mathit{callees}}
\newcommand{\nul}{\mathsf{null}}
\newcommand{\tok}{\bullet}
\newcommand{\ttrue}{\mathit{tt}}
\newcommand{\tfalse}{\mathit{ff}}
\newcommand{\pof}[1]{[\![#1]\!]}
\newcommand{\up}[1]{\lceil #1 \rceil}
\newcommand{\cfg}[1]{\langle #1 \rangle}

An optimizing compiler may reorder, merge and remove operations on memory only
when nothing in between can interfere with them. Within a method it can usually
tell. A call it has not inlined is another matter: Graal~\cite{duboscq2013ir},
like most compilers, takes a call to read and write all of memory, so it
forwards no load across a call and removes no store before one. The running
example of this chapter is a method of the program \texttt{Aliasing}, one of
those the thesis evaluates (\cref{chap:evaluation}):

\begin{lstlisting}[language=Java]
static int forwardAcrossCallOnOtherObject(Counter counter, Counter other, Link list) {
    counter.value = 5;
    countInto(other, list);   // writes Counter.value, but only of other
    return counter.value;
}
\end{lstlisting}

The call writes the field \texttt{value} of the counter it is given, and the
program only ever gives it a counter other than the one it passes this method
as \texttt{counter}. The load after the call therefore reads 5. Graal keeps the
load, because in its model the call may write anything.

This chapter makes the effects of each operation explicit, as values. Every
method takes extra parameters and gives back extra results, \emph{tokens}: one
for each part of memory, a \emph{partition}, that it or anything it calls may
touch. Every operation on memory takes the token of its partition and gives
back a new one, which the next operation on that partition takes. The order of
the operations on memory becomes data flow. In a program in which every value is
used exactly once, the tokens of a partition form a chain, and two operations
are ordered only when a chain connects them. In the example, with a partition
for each field of each set of objects the points-to analysis tells apart, the
call takes the token of the counters the program passes as \texttt{other}, but
not the token of the one it passes as \texttt{counter}. The load's token comes
straight from the store, so the load can be replaced by the 5.

\Cref{sec:ee-language} defines a small language with tokens, the core of
graal.ct, the dialect of Cthulhu in which the implementation represents Graal
IR, and \cref{sec:ee-semantics} its semantics. \Cref{sec:ee-partitions}
defines partitions and the four partitionings the implementation offers, and
\cref{sec:ee-signatures} the signatures that say which tokens a method takes.
\Cref{sec:ee-guarantee} proves what the tokens guarantee: an operation leaves
a partition alone unless it takes its token, and two operations that share no
token can be swapped. \Cref{sec:ee-license} proves which rewrites the tokens
license: forwarding a load and removing a store, in a model in which an
exception ends the program and in Java, and that a finer partitioning licenses
every rewrite a coarser one does. The next three sections give the conditions
under which an effect may be left out: initializing a new object
(\cref{sec:ee-fresh}) and synchronization (\cref{sec:ee-concurrency}), and what
the same facts guarantee when they go back to Graal as kill sets
(\cref{sec:ee-kills}). \Cref{sec:ee-summary} lists every result with what it
assumes and what relies on it.

Each result was also held against the implementation. Where the two disagree,
the section says so, and \cref{sec:ee-summary} collects what that turned up.

%----------------------------------------------------------------------------
\section{A language with tokens}\label{sec:ee-language}

The implementation translates each method of Graal IR into graal.ct
(\cref{chap:conversion}). We reason about a small core of graal.ct, which keeps
everything that matters for memory. \Cref{fig:ee-syntax} gives its syntax.

\begin{figure}[t]
\[
\begin{array}{@{}r@{\;\;}c@{\;\;}l@{\qquad}l@{}}
\tau & ::= & \mathsf{int} \mid \mathsf{bool} \mid \mathsf{ref} & \text{values}\\
     & \mid & \Tok{P} \qquad (P \in \Pi) & \text{the token of partition } P\\[6pt]
\iota & ::= & \bar x \coloneqq \mathit{op}(\bar y) & \text{an operation on values}\\
 & \mid & y_1, y_2 \coloneqq \dup\ x \;\mid\; \drop\ x & \text{copy, discard a value}\\
 & \mid & h' \coloneqq \move\ h & \text{pass a token on}\\
 & \mid & x, h' \coloneqq \getf{f}{P}(h, y) & \text{load field } f \text{ of } y\\
 & \mid & h' \coloneqq \setf{f}{P}(h, y, z) & \text{store } z \text{ into field } f \text{ of } y\\
 & \mid & h' \coloneqq \initf{f}{P}(h, y, z) \;\mid\; \initf{f}{}(y, z) & \text{give field } f \text{ of a new } y \text{ its first value}\\
 & \mid & x, h' \coloneqq \newc{C}(h) & \text{allocate an object of class } C\\
 & \mid & \bar h' \coloneqq \guard(\bar h, y) & \text{trap when } y \text{ holds}\\
 & \mid & \bar x, \bar h' \coloneqq m(\bar y, \bar h) & \text{call method } m\\[6pt]
\kappa & ::= & \jump\ b(\bar y) & \text{go on in block } b\\
 & \mid & \kif{y}{b_1}{b_2}\,(\bar y) & \text{branch}\\
 & \mid & \kret(\bar y) \;\mid\; \ktrap(\bar h) & \text{leave the method}\\[6pt]
B & ::= & b(\bar x : \bar \tau) = \iota_1 \ldots \iota_n;\ \kappa & \text{a block}
\end{array}
\]
\caption{The core of graal.ct. Variables $x, y, z$ hold values and $h$
tokens; $\bar x$ is a sequence of them. A method is a set of blocks, one of
which is its entry.}
\label{fig:ee-syntax}
\end{figure}

A program consists of methods. A method $m$ consists of blocks, one of which
is its entry block $b_m$. It takes values of types $\bar\tau_m$ and gives back
values of types $\bar\tau'_m$, and it takes and gives back one token of each
partition in a set $T_m$, defined in \cref{sec:ee-signatures}. A block
$b(\bar x : \bar\tau) = \iota_1 \ldots \iota_n;\ \kappa$ binds its parameters
$\bar x$, runs its instructions in order, and ends with a terminator $\kappa$:
a jump to another block of the method, a branch to one of two, or leaving the
method, with $\kret$ or with $\ktrap$. Blocks interact only through their
parameters, and methods only through calls and memory.

This is the shape of Graal IR once it is scheduled: a basic block becomes a
block, a $\varphi$-function becomes a parameter of its block, and each edge
passes the $\varphi$'s input for that edge as an argument. That a program in
static single assignment form can be read this way, as a set of mutually
recursive functions that call each other in tail position, is well
known~\cite{kelsey1995cps,appel1998ssa}.

Each memory instruction is annotated with a partition $P$. A load
$x, h' \coloneqq \getf{f}{P}(h, y)$ takes the token $h$ of $P$, reads field
$f$ of the object $y$ into $x$, and gives back the next token of $P$ as $h'$.
A store gives back only the token. An initialization gives a field of an object
that was just allocated its first value; it takes a token only if its method
takes the token of that field's partition for another reason
(\cref{sec:ee-fresh}). An allocation takes the token of the rest of memory,
$\Rest$, which holds the allocator's state. A guard traps when its condition
holds, and takes and gives back a token of every partition of its method. A
call takes and gives back the tokens of the callee's partitions.

\subsection{Single use}

\begin{definition}[Static single use]\label{def:ee-ssu}
A method is in \emph{static single use} (SSU) form if each of its blocks
satisfies the following:
\begin{enumerate}
\item every variable is bound once, as a parameter of the block or as a result
  of one of its instructions;
\item every variable bound in the block is used exactly once after it is bound,
  by an instruction or by the terminator;
\item $\dup$ and $\drop$ are only applied to values, never to tokens.
\end{enumerate}
\end{definition}

Static single assignment~\cite{cytron1991ssa} asks only for the first condition,
and lets a value be used any number of times. Static single use asks that every
use have a value of its own: a value two instructions need is first copied
with $\dup$, and a value nobody needs is discarded with $\drop$. What this costs
is the copies. What it buys is that the lifetime of a value is explicit, from
its definition to its one use.

The third condition is what makes tokens special. A variable that may be
neither copied nor discarded is \emph{linear}~\cite{girard1987linear,
wadler1990linear}: it is used exactly once, so at every point of a block
exactly one version of each token exists. Threading the state of the world
through a program as a linear value is an old idea. Haskell's state
threads~\cite{launchbury1994state} and Clean's uniqueness
types~\cite{barendsen1996uniqueness} use it to make input, output and mutable
arrays purely functional. Here the state is split, one token per partition,
and the split is what lets operations on different partitions happen in any
order.

\subsection{Types}

\begin{figure}[t]
\centering
\begin{tabular}{@{}lll@{}}
\toprule
Instruction & Takes & Gives \\
\midrule
$\bar x \coloneqq \mathit{op}(\bar y)$ & the operation's values & its values \\
$y_1, y_2 \coloneqq \dup\ x$; $\drop\ x$ & a value of type $\tau \neq \Tok{P}$ & two of type $\tau$; nothing \\
$h' \coloneqq \move\ h$ & $\Tok{P}$ & $\Tok{P}$ \\
$x, h' \coloneqq \getf{f}{P}(h, y)$ & $\Tok{P}$, $\mathsf{ref}$ & $\tau_f$, $\Tok{P}$ \\
$h' \coloneqq \setf{f}{P}(h, y, z)$, $h' \coloneqq \initf{f}{P}(h, y, z)$ & $\Tok{P}$, $\mathsf{ref}$, $\tau_f$ & $\Tok{P}$ \\
$\initf{f}{}(y, z)$ & $\mathsf{ref}$, $\tau_f$ & nothing \\
$x, h' \coloneqq \newc{C}(h)$ & $\Tok{\Rest}$ & $\mathsf{ref}$, $\Tok{\Rest}$ \\
$\bar h' \coloneqq \guard(\bar h, y)$ & $\Tok{S}$, $\mathsf{bool}$ & $\Tok{S}$ \\
$\bar x, \bar h' \coloneqq m(\bar y, \bar h)$ & $\bar\tau_m$, $\Tok{T_m}$ & $\bar\tau'_m$, $\Tok{T_m}$ \\
\bottomrule
\end{tabular}
\caption{The types of instructions. $\tau_f$ is the type of field $f$, and
$\Tok{S}$ for a set $S$ of partitions is one token of each, in a fixed order.
Every instruction gives back exactly the tokens it takes.}
\label{fig:ee-types}
\end{figure}

\begin{definition}[Well-typed method]\label{def:ee-typed}
A method $m$ with tokens $T_m$ is \emph{well typed} if every block of $m$
takes, besides its values, one token of each partition in $T_m$; every
instruction is used at its type in \cref{fig:ee-types}, with guards over
$S = T_m$; the arguments of a jump or a branch have the types of the
parameters of the blocks it goes to, and the two blocks of a branch take the
same parameters; and $\kret$ gives back values of types $\bar\tau'_m$ and a
token of each partition in $T_m$, which $\ktrap$ takes as well.
\end{definition}

\begin{lemma}[Chains]\label{lem:ee-chains}
Let $b$ be a block of a well-typed method $m$ in SSU form, and $P \in T_m$. The
instructions of $b$ that take a token of $P$, in the order in which they appear,
are $\iota_{k_1}, \ldots, \iota_{k_r}$, where $\iota_{k_1}$ takes the parameter
of $b$ of type $\Tok{P}$, each $\iota_{k_{i+1}}$ takes the token of $P$ that
$\iota_{k_i}$ gives, and the terminator takes the one $\iota_{k_r}$ gives, or
the parameter if $r = 0$.
\end{lemma}

\begin{proof}
We show that after each instruction exactly one variable of type $\Tok{P}$ is
bound and not yet used. Before the first instruction, it is the parameter.
An instruction that takes a token of $P$ takes exactly one
(\cref{fig:ee-types}), and it must be that variable: every other variable of
type $\Tok{P}$ bound so far has been used already, and none can be used twice.
The instruction gives back exactly one new token of $P$. An instruction that
takes no token of $P$ gives none. After the last instruction, the one
remaining token must still be used, and only the terminator is left.
\end{proof}

So in every block the operations on one partition are totally ordered, and the
order is the token's data flow. Nothing orders two instructions that take the
tokens of different partitions, other than the values one gives the other.

\subsection{Branches in Cthulhu}\label{sec:ee-branches}

Cthulhu has no branch terminator: every block ends with a call. A two-way
branch is built from closures, as \cref{fig:ee-branch} shows on the running
example. $\mathsf{opt}(c, k)$ is the closure $k$ when $c$ holds and $\bot$
otherwise, and $\mathsf{join}(a_1, a_2, g)$ is a closure that calls the
\emph{frame} $g$ with both. The frame forks each token in two, gives one half to
each arm, and joins what the two arms give back, where $\bot$ is the unit of
$\mathsf{join}$. In every execution exactly one of the arms is not $\bot$, so of
the two halves of a token, only one ever reaches an operation: $\mathsf{fork}$
is not $\dup$. In terms of linear logic, the two halves form the additive
conjunction of the two arms, not the multiplicative one: a block may use either,
never both. The core language replaces the encoding with $\mathsf{if}$, which
behaves the same.

\begin{figure}[t]
\begin{lstlisting}[language=ct]
run = λ p0 p1 p2 h → r hout                          ; ref × ref × ref × heap → int × heap
(
    ref nil? p0 → v6                                     ; counter == null
    Aliasing_forwardAcrossCallOnOtherObject b1 → kt      ; b1 throws
    Aliasing_forwardAcrossCallOnOtherObject b2 → kf
    Aliasing_forwardAcrossCallOnOtherObject frame_0 → frame
    bool not v6 → cn
    f_rrrh_ih opt v6 kt → at
    f_rrrh_ih opt cn kf → af
    f_rrrh_ih join at af frame → k
    f_rrrh_ih call k p0 p1 p2 h → r hout
)
frame_0 = λ A B x0 x1 x2 h → r hout
(
    heap fork h → h_a h_b
    f_rrrh_ih call A x0 x1 x2 h_a → r_a hout_a
    f_rrrh_ih call B x0 x1 x2 h_b → r_b hout_b
    int join r_a r_b → r
    heap join hout_a hout_b → hout
)
\end{lstlisting}
\caption{The null check at the start of the running example, with one token
for all of memory, as the implementation prints it but without the copies of
values, which $\dup$ makes. \texttt{f\_rrrh\_ih} is the type of a closure that
takes three references and a token and gives back an \texttt{int} and a token.}
\label{fig:ee-branch}
\end{figure}

\subsection{The rest of graal.ct}

graal.ct has more than the core. It has numbers of four sizes and their
operations, constants of the program, and static fields, which it accesses
through one object per class. It has arrays, whose elements are locations
$(o, i)$ and whose accesses $\mathsf{aget}_K$ and $\mathsf{aset}_K$ are annotated
like field accesses: with the partition of the element kind $K$, or of the
elements of kind $K$ of a set of arrays. And it has three kinds of call the core
lacks:
\begin{itemize}
\item a node of Graal IR that touches memory in a way graal.ct has no operation
  for, such as an array copy or the initialization of a class, becomes a call
  of an external method that takes every token of its method;
\item a call with several targets becomes a call of an external method that
  takes the tokens of all of them;
\item a call the analysis could not resolve takes every token of the program.
\end{itemize}
All of them are memory instructions or calls in the sense of this chapter: each
takes the tokens of every partition its execution may touch, which is all the
results below rely on.

%----------------------------------------------------------------------------
\section{Semantics}\label{sec:ee-semantics}

\begin{definition}[Heaps]\label{def:ee-heaps}
Objects are natural numbers, $\Obj = \mathbb{N}$, and values are
$\Val = \mathbb{Z} \cup \{\ttrue, \tfalse\} \cup \Obj \cup \{\nul, \tok\}$,
where $\tok$ is the only value a token ever has. A \emph{location} is a field of
an object, or the allocator: $\Loc = (\Obj \times F) \cup \{\alloc\}$, where $F$
is the set of fields, including two header fields, $\mathsf{cls}$ and
$\mathsf{site}$, which record an object's class and where it was allocated,
and which no instruction writes. A \emph{heap} is a finite partial function
$\sigma$ from $\Loc$ to $\Val$ in which $\sigma(\alloc)$ is the number of
objects allocated so far. An object $o$ \emph{exists} in $\sigma$ if
$o < \sigma(\alloc)$, and then $(o, f) \in \operatorname{dom} \sigma$ for each
field $f$ of its class and for the header fields. An \emph{environment} $\rho$
maps variables to values.
\end{definition}

\begin{figure}[tbp]
\begin{mathpar}
\inferrule*[left=Op]{ }{\cfg{\bar x \coloneqq \mathit{op}(\bar y), \rho, \sigma}
  \Downarrow \cfg{\rho[\bar x \mapsto \pof{\mathit{op}}(\rho(\bar y))], \sigma}}

\inferrule*[left=Get]{\rho(y) = o \neq \nul}{\cfg{x, h' \coloneqq \getf{f}{P}(h, y), \rho, \sigma}
  \Downarrow \cfg{\rho[x \mapsto \sigma(o, f),\ h' \mapsto \tok], \sigma}}

\inferrule*[left=Set]{\rho(y) = o \neq \nul}{\cfg{h' \coloneqq \setf{f}{P}(h, y, z), \rho, \sigma}
  \Downarrow \cfg{\rho[h' \mapsto \tok], \sigma[(o, f) \mapsto \rho(z)]}}

\inferrule*[left=Null]{\rho(y) = \nul}{\cfg{\iota, \rho, \sigma} \Downarrow \Trap}

\inferrule*[left=New]{o = \sigma(\alloc)}{\cfg{x, h' \coloneqq \newc{C}(h), \rho, \sigma}
  \Downarrow \cfg{\rho[x \mapsto o,\ h' \mapsto \tok],\ \sigma[\alloc \mapsto o + 1][(o, f) \mapsto 0_f]_{f}}}

\inferrule*[left=Guard]{\rho(y) = \tfalse}{\cfg{\bar h' \coloneqq \guard(\bar h, y), \rho, \sigma}
  \Downarrow \cfg{\rho[\bar h' \mapsto \tok], \sigma}}

\inferrule*[left=Fail]{\rho(y) = \ttrue}{\cfg{\bar h' \coloneqq \guard(\bar h, y), \rho, \sigma} \Downarrow \Trap}

\inferrule*[left=Call]{\cfg{m, \rho(\bar y), \sigma} \Downarrow \Ret(\bar v, \sigma')}{\cfg{\bar x, \bar h' \coloneqq m(\bar y, \bar h), \rho, \sigma}
  \Downarrow \cfg{\rho[\bar x \mapsto \bar v,\ \bar h' \mapsto \tok], \sigma'}}

\inferrule*[left=Throw]{\cfg{m, \rho(\bar y), \sigma} \Downarrow \Trap}{\cfg{\bar x, \bar h' \coloneqq m(\bar y, \bar h), \rho, \sigma} \Downarrow \Trap}
\\
\inferrule*[left=Seq]{\cfg{\iota, \rho, \sigma} \Downarrow \cfg{\rho', \sigma'} \\ \cfg{\bar\iota; \kappa, \rho', \sigma'} \Downarrow r}{\cfg{\iota\,\bar\iota; \kappa, \rho, \sigma} \Downarrow r}

\inferrule*[left=Stop]{\cfg{\iota, \rho, \sigma} \Downarrow \Trap}{\cfg{\iota\,\bar\iota; \kappa, \rho, \sigma} \Downarrow \Trap}

\inferrule*[left=Jump]{\cfg{b, \rho(\bar y), \sigma} \Downarrow r}{\cfg{\jump\ b(\bar y), \rho, \sigma} \Downarrow r}

\inferrule*[left=If]{\rho(y) = \ttrue \\ \cfg{b_1, \rho(\bar y), \sigma} \Downarrow r}{\cfg{\kif{y}{b_1}{b_2}\,(\bar y), \rho, \sigma} \Downarrow r}

\inferrule*[left=Return]{ }{\cfg{\kret(\bar y), \rho, \sigma} \Downarrow \Ret(\rho(\bar y), \sigma)}

\inferrule*[left=Trap]{ }{\cfg{\ktrap(\bar h), \rho, \sigma} \Downarrow \Trap}

\inferrule*[left=Block]{b(\bar x : \bar\tau) = \bar\iota; \kappa \\ \cfg{\bar\iota; \kappa, [\bar x \mapsto \bar v], \sigma} \Downarrow r}{\cfg{b, \bar v, \sigma} \Downarrow r}
\end{mathpar}
\caption{Evaluation. An outcome $r$ is $\Ret(\bar v, \sigma)$ or $\Trap$. An
initialization, with a token or without, evaluates as \textsc{Set}; $\dup$,
$\drop$ and $\move$ only change the environment; \textsc{Null} applies to a
load, a store and an initialization through $\nul$; $\mathsf{if}$ with
$\rho(y) = \tfalse$ goes to $b_2$; a method is invoked as its entry block,
with $\tok$ for each token, $\cfg{m, \bar v, \sigma} \Downarrow r$ if
$\cfg{b_m, \bar v\,\bar\tok, \sigma} \Downarrow r$, and $\Ret$ gives back
only the values. $0_f$ is the default value of field $f$, and
\textsc{New} also sets the header fields.}
\label{fig:ee-semantics}
\end{figure}

\Cref{fig:ee-semantics} gives the semantics. The relation is defined
inductively, so an invocation that does not terminate has no outcome.
Operations on values are total. That is what Graal IR gives us: Native Image
checks a divisor, an array index, and the receiver of an access explicitly,
before the operation itself, with a branch to a block that throws, which the
translation makes a block that traps. Besides $\ktrap$, only guards and
accesses through $\nul$ trap, and the latter never execute: every access
comes after the check of its receiver, or Graal knows the receiver is not
$\nul$.

\begin{lemma}[Determinism]\label{lem:ee-determinism}
For every method $m$, values $\bar v$ and heap $\sigma$ there is at most one
outcome $r$ with $\cfg{m, \bar v, \sigma} \Downarrow r$.
\end{lemma}

\begin{proof}
By induction on the derivation: the premises of every rule determine its
conclusion, and an allocation takes the object $\sigma(\alloc)$.
\end{proof}

\begin{lemma}[Tokens carry nothing]\label{lem:ee-erasure}
Let two programs differ only in their tokens: in which token variables their
instructions and terminators take and give, and in moves of tokens. Then every
invocation has the same outcome in both.
\end{lemma}

\begin{proof}
Every token variable holds $\tok$, and no rule looks at the value of a token. By
induction on the derivation.
\end{proof}

The lemma has a consequence the rest of the chapter relies on: tokens add no
behaviour to a program. A program annotated with the tokens of one partitioning
or of another (\cref{sec:ee-signatures}) is the same program, with the same
outcomes. What a partitioning changes is which rewrites its tokens license. In
particular, the tokens order operations, but they do not hold the heap: an
interpreter that gave each token a copy of its partition would be wrong for the
programs of this thesis, which leave the tokens out of some initializations
(\cref{sec:ee-fresh}).

\begin{definition}[Runs]\label{def:ee-runs}
The \emph{run} of an invocation $\cfg{m, \bar v, \sigma}$ is the sequence of
\emph{steps} $s_1, s_2, \ldots$ that its evaluation executes in the blocks of
$m$ itself, through jumps and branches: step $s_k$ executes an instruction
$\iota_k$ in environment $\rho_k$ and heap $\sigma_k$, and leaves the heap
$\sigma_{k+1}$. A call is one step. The run ends with the terminator that
gives the outcome; the run of an invocation that does not terminate is
infinite, or ends with a call that does not terminate. A step \emph{takes} $P$
if its instruction takes a token of $P$. A location \emph{exists} at step $k$
if its object exists in $\sigma_k$; objects are never removed, so it exists at
every later step too.
\end{definition}

The results below are about runs of the program: the invocation of its
\texttt{main} method from the initial heap, and every invocation it makes.
That is the only place where the facts of the points-to analysis hold.

%----------------------------------------------------------------------------
\section{Partitions}\label{sec:ee-partitions}

A partition is a set of locations, and a partitioning divides all of memory
into disjoint partitions. Which partitions are useful depends on what is known
about the objects an operation may touch, and that knowledge comes from a
points-to analysis.

\subsection{Abstract objects}

A points-to analysis describes the objects a reference may point to by
\emph{abstract objects}. Native Image's analysis~\cite{wimmer2019initialize}
describes an object by its class by default, and with allocation sites by the
place in the program where it was allocated. We write $\alpha(o) \in \AObj$ for
the abstract object of $o$. It is determined when $o$ is allocated, by the
header fields $\mathsf{cls}$ and $\mathsf{site}$, and never changes.

\begin{assumption}[Sound points-to]\label{as:ee-points-to}
Each variable $y$ that an access uses as its receiver has a set
$\pts(y) \subseteq \AObj$, or $\pts(y) = \top$ when the analysis cannot say.
In every run of the program, whenever an access through $y$ executes with
$\rho(y) = o \neq \nul$, then $\pts(y) = \top$ or $\alpha(o) \in \pts(y)$.
\end{assumption}

\subsection{Partitionings}

\begin{definition}[Partitioning]\label{def:ee-partitioning}
A \emph{partitioning} is a set $\Pi$ of nonempty, pairwise disjoint sets of
locations, \emph{partitions}, whose union is $\Loc$, such that the partition of
a field $(o, f)$ depends only on $f$ and $\alpha(o)$. We write $\pof{\ell}$ for
the partition that contains $\ell$, and $\Rest$ for the one that contains
$\alloc$.
\end{definition}

Since $\alpha(o)$ never changes, a location never changes its partition.

\begin{definition}[Labelling]\label{def:ee-labelling}
A \emph{labelling} gives each access $s$ of a program, a load, a store or an
initialization of field $f$, a partition $\lab(s) \in \Pi$. It is \emph{sound}
if in every run of the program every execution of $s$ accesses a location of
$\lab(s)$.
\end{definition}

The four partitionings of the implementation differ in which locations must
share a partition. Let $F_p$ be the fields of the program's classes; every other
field is the platform's, from the JDK or from Native Image. For $f \in F_p$, let
$\mathit{Known}(f)$ be the abstract objects that the accesses of $f$ through
references the analysis can describe may touch: the union of $\pts(y)$ over
those accesses. An \emph{abstract field} is a pair $(f, a)$ with
$a \in \mathit{Known}(f)$, or $(f, ?)$, which stands for every object no such
access names. The abstract field of a location $(o, f)$ is $(f, \alpha(o))$ if
$\alpha(o) \in \mathit{Known}(f)$, and $(f, ?)$ otherwise. An access $s$ of
$f \in F_p$ through $y$ \emph{may touch} the abstract fields
$E(s) = \{f\} \times \pts(y)$, or every abstract field of $f$ if
$\pts(y) = \top$.

\begin{definition}[The four partitionings]\label{def:ee-four}
For an equivalence $\equiv$ on abstract fields, let $\Pi_\equiv$ be the
partitioning in which two fields $(o, f)$ and $(o', f')$ with $f, f' \in F_p$
are in one partition if and only if their abstract fields are related by
$\equiv$, and the platform's fields, the header fields and $\alloc$ form one
more partition, $\Rest$. Let
\begin{itemize}
\item $G_{\mathrm{acc}}$ relate any two abstract fields that one access may touch,
\item $G_{\mathrm{obj}}$ relate $(f, a)$ and $(g, a)$ for every abstract object
  $a$ and fields $f, g$ that accesses name it with: the fields of one object,
\item $G_{\mathrm{fld}}$ relate any two abstract fields of one field $f$: one
  field of all objects,
\end{itemize}
and write $G^{\ast}$ for the least equivalence that contains $G$. Then
\emph{object-field} is $\Pi_\equiv$ for ${\equiv} = G_{\mathrm{acc}}^{\ast}$,
\emph{object} for $(G_{\mathrm{acc}} \cup G_{\mathrm{obj}})^{\ast}$, and \emph{field}
for $(G_{\mathrm{acc}} \cup G_{\mathrm{fld}})^{\ast} = G_{\mathrm{fld}}^{\ast}$; \emph{none}
is $\{\Loc\}$. The label of an access $s$ of $f \in F_p$ is the partition of the
abstract fields in $E(s)$, the label of an initialization is the partition of
the abstract field of the location it initializes, and the label of an access
of a platform field is $\Rest$.
\end{definition}

\begin{table}[t]
\centering
\begin{tabular}{@{}lll@{}}
\toprule
Partitioning & A partition of the program's fields is & Tells apart \\
\midrule
none & all of memory, in one & nothing \\
field & one field of all objects & fields, as Graal's own locations do \\
object & every field of a set of objects & objects one access never mixes \\
object-field & one field of a set of objects & both \\
\bottomrule
\end{tabular}
\caption{The four partitionings. Under every one but none, the platform's
fields and the allocator are one more partition, $\Rest$.}
\label{tab:ee-partitionings}
\end{table}

\Cref{tab:ee-partitionings} sums them up. The label is well defined because
$G_{\mathrm{acc}}$ relates all abstract fields in $E(s)$, so they lie in one
partition. The implementation computes object and object-field by union-find
over the abstract fields the facts list~\cite{steensgaard1996points}: it joins
the abstract fields of each access, and for a field that some access touches
through a reference the analysis can't describe, it joins $(f, ?)$ with every
object the other accesses of that field may touch. That gives the same
partitions as the definition for every location an access may touch.

\begin{proposition}[Sound labelling]\label{prop:ee-sound-labels}
Under \cref{as:ee-points-to}, each of the four is a partitioning, and its
labelling is sound.
\end{proposition}

\begin{proof}
The partitions are the classes of an equivalence, and $\Rest$, so they are
disjoint and cover $\Loc$, and the partition of $(o, f)$ depends only on
$(f, \alpha(o))$. Let an execution of an access $s$ of field $f$ through $y$
access $(o, f)$. If $f$ is a platform field, $\lab(s) = \Rest \ni (o, f)$.
Otherwise the abstract field of $(o, f)$ is in $E(s)$: if $\pts(y) = \top$,
$E(s)$ holds every abstract field of $f$, and otherwise
$\alpha(o) \in \pts(y) \subseteq \mathit{Known}(f)$ by \cref{as:ee-points-to},
so the abstract field is $(f, \alpha(o)) \in E(s)$. Hence the partition of
$(o, f)$ is that of $E(s)$, which is $\lab(s)$. An initialization is labelled
by its location's partition.
\end{proof}

\subsection{Finer and coarser}

\begin{definition}[Refinement]\label{def:ee-refinement}
A partitioning $\Pi'$ \emph{refines} $\Pi$, written $\Pi' \sqsubseteq \Pi$, if
every partition of $\Pi'$ is a subset of a partition of $\Pi$. For
$P' \in \Pi'$ we write $\up{P'}$ for that partition of $\Pi$.
\end{definition}

\begin{proposition}[The order of the four]\label{prop:ee-order}
object-field $\sqsubseteq$ field $\sqsubseteq$ none, and object-field
$\sqsubseteq$ object $\sqsubseteq$ none, while field and object are in general
incomparable. Whenever $\Pi' \sqsubseteq \Pi$ are two of them, every access
$s$ has $\lab_{\Pi'}(s) \subseteq \lab_\Pi(s)$, that is,
$\up{\lab_{\Pi'}(s)} = \lab_\Pi(s)$.
\end{proposition}

\begin{proof}
If $G \subseteq G'$, then every class of $G^{\ast}$ lies within a class of
$G'^{\ast}$, and $G_{\mathrm{acc}}$ is contained in both of the larger relations.
$\Rest$ is the same set in all four, and everything lies within the one
partition of none. The label of $s$ under $\Pi'$ is the partition that contains
$E(s)$, which lies within the partition of $\Pi$ that contains $E(s)$. For
incomparability, see \cref{ex:ee-incomparable}.
\end{proof}

\begin{example}[Field and object are incomparable]\label{ex:ee-incomparable}
The method \texttt{deadStoreOtherField} of \texttt{Examples} stores to
\texttt{value}, then to \texttt{other}, then to \texttt{value} again, all on one
counter. With a partition per field, the store to \texttt{other} takes another
token, so the first store's token goes straight to the third store, which
overwrites it: the first store is dead (\cref{sec:ee-license}). With a
partition per set of objects, all three take the counter's token, and the store
to \texttt{other} stands in between. Conversely, in the running example
\texttt{counter} and \texttt{other} are told apart only by the objects. The
finest of the four, object-field, tells apart what either of them does.
\end{example}

%----------------------------------------------------------------------------
\section{Signatures}\label{sec:ee-signatures}

A method's signature says which partitions it, and everything it calls, may
read and write. It is made of facts about each method, which come from the
points-to analysis and from Graal IR (\cref{chap:conversion}): the accesses it
makes, with their receivers; the methods each of its calls may invoke; whether
it allocates; and every other node of Graal IR that touches memory, with the
locations Graal says the node touches.

\begin{assumption}[Sound facts]\label{as:ee-facts}
\begin{enumerate}
\item For every call $s$, the set $\callees(s)$ contains every method the call
  invokes in a run of the program.
\item Every other node that touches memory, other than an access, an allocation
  or a call, touches only locations in the partitions of fields and array kinds
  the facts list among its method's reads and writes, or makes its method
  opaque.
\end{enumerate}
A method is \emph{opaque} if it contains a node that may touch a location the
facts cannot name (a location Graal does not name, a lock, an atomic update, a
volatile access, a transition to native code), a call the analysis could not
resolve, or a call of an opaque method.
\end{assumption}

\begin{definition}[Signature]\label{def:ee-signature}
Let $\Pi$ be a partitioning with a labelling. The \emph{signature} of a method
$m$ is the pair $(R(m), W(m))$ of sets of partitions that is the least solution
of the constraints
\begin{align*}
R(m) &\supseteq \{\lab(s) \mid s \text{ a load in } m\},\\
W(m) &\supseteq \{\lab(s) \mid s \text{ a store in } m\} \cup \{\Rest \mid m \text{ allocates}\},\\
R(m) &\supseteq R(m') \text{ and } W(m) \supseteq W(m') \quad \text{for every } m' \in \callees(s) \text{ of a call } s \text{ in } m,
\end{align*}
together with the partitions of the locations the other nodes of $m$ read and
write. The \emph{tokens} of $m$ are $T_m = R(m) \cup W(m)$, or $T_m = \Pi$ if $m$
is opaque.
\end{definition}

The constraints are monotone, and their values form a finite lattice, so a least
solution exists and iterating them from the empty sets reaches
it~\cite{tarski1955lattice}. The implementation computes it so
(\cref{chap:conversion}). Initializations are not among the stores: the facts
leave them out, for the reason \cref{sec:ee-fresh} gives.

\begin{definition}[Token form]\label{def:ee-token-form}
The \emph{token form} of a program under $\Pi$ gives its instructions these
tokens:
\begin{itemize}
\item every block of a method $m$ takes the tokens of $T_m$;
\item a load and a store take the token of their label;
\item an initialization takes the token of its label if that is in $T_m$, and no
  token otherwise;
\item an allocation takes $\Rest$;
\item a call takes the tokens $T_{m'}$ of its callee $m'$, the union of $T_{m'}$
  over its callees if it has several, and all of $\Pi$ if the analysis could not
  resolve it;
\item a guard, and a node of the second kind in \cref{as:ee-facts}, takes the
  tokens $T_m$ of its method.
\end{itemize}
\end{definition}

\begin{lemma}[The token form is well typed]\label{lem:ee-token-form-typed}
The token form of every method is well typed.
\end{lemma}

\begin{proof}
Every token an instruction of $m$ takes is in $T_m$. For loads, stores and
allocations, this is the first line of \cref{def:ee-signature}. A callee's
tokens are in $T_m$ by the second line, and if the callee is opaque, so is $m$,
and $T_m = \Pi$. An unresolved call makes $m$ opaque too. The rest take $T_m$
by definition.
\end{proof}

\begin{example}[The running example]\label{ex:ee-running-tokens}
The allocation-site analysis gives the two counters \texttt{main} passes to
\texttt{forwardAcrossCallOnOtherObject} the abstract objects
\texttt{Counter@main:199}, for \texttt{counter}, and \texttt{Counter@main:206},
for \texttt{other}. No access ever touches both, and \texttt{countInto} is only
ever given the second. Under object-field, \texttt{Counter.value} therefore has
two partitions, \texttt{value@main:199} and \texttt{value@main:206}, and
\begin{align*}
W(\texttt{countInto}) &= \{\texttt{value@main:206}\},\\
R(\texttt{countInto}) &= \{\texttt{value@main:206},\ \texttt{Link.next}\}.
\end{align*}
\Cref{fig:ee-running-tokens} shows the two blocks of the method that hold the
store, the call and the load. The call takes the token of \texttt{other}'s
counters and that of \texttt{Link.next}, and passes the token of
\texttt{counter}'s straight from the store to the load. Under field the call
takes \texttt{h\_value} and \texttt{h\_next}, and with one partition it takes
\texttt{h}: in both the load's token comes from the call.
\end{example}

\begin{figure}[tbp]
\begin{lstlisting}[language=ct]
b2 = λ p0 p1 p2 h_c h_o h_n → r h_cout h_oout h_nout
(
    int cons_5 → v11
    heap set_0 h_c p0 v11 → h_c1                         ; counter.value = 5
    Aliasing_countInto run → k13
    f_rrhh_hh call k13 p1 p2 h_o h_n → h_o1 h_n1         ; countInto(other, list)
    Aliasing_forwardAcrossCallOnOtherObject b3 → k3
    f_rhhh_ihhh call k3 p0 h_c1 h_o1 h_n1 → r h_cout h_oout h_nout
)
b3 = λ p0 h_c h_o h_n → r h_cout h_oout h_nout
(
    heap get_0 h_c p0 → v24 h_c1                         ; return counter.value
    int move v24 → r
    heap move h_c1 → h_cout
    heap move h_o → h_oout
    heap move h_n → h_nout
)
\end{lstlisting}
\caption{The running example under object-field, as the implementation prints
it but without the copies of values, and with shorter names for the tokens:
\texttt{h\_c} for \texttt{Counter.value@main:199}, \texttt{h\_o} for
\texttt{Counter.value@main:206}, and \texttt{h\_n} for \texttt{Link.next}.
Field 0 is \texttt{Counter.value}.}
\label{fig:ee-running-tokens}
\end{figure}

\subsection{What a signature covers}

\begin{definition}[Footprint]\label{def:ee-footprint}
An invocation \emph{reads} a location if one of the steps of its run is a load
of it, an allocation (which reads $\alloc$), or a call whose invocation reads
it, and it \emph{writes} a location if one of its steps is a store or an
initialization of it, an allocation (which writes $\alloc$), or a call whose
invocation writes it. The \emph{footprint} of an invocation
$\cfg{m, \bar v, \sigma}$ is made of the locations that exist in $\sigma$, when
it starts, that it reads and writes.
\end{definition}

\begin{theorem}[Signatures]\label{thm:ee-signatures}
Assume \cref{as:ee-points-to,as:ee-facts}, and let $\Pi$ be one of the four
partitionings. For every invocation $\cfg{m, \bar v, \sigma}$ in a run of the
program, every location in its footprint lies in a partition of $T_m$, and if
$m$ is not opaque, every location of its footprint that it writes lies in a
partition of $W(m)$.
\end{theorem}

\begin{proof}
If $m$ is opaque, $T_m = \Pi$ covers every location, so let $m$ not be opaque.
We proceed by induction on the derivation of the invocation's outcome. Let
$\ell$ exist in $\sigma$, and let a step of the run write $\ell$; reads are
similar, with $R(m)$ in place of $W(m)$ and $\Rest \in W(m)$ for the reads of an
allocation.
\begin{itemize}
\item A store: by \cref{prop:ee-sound-labels}, $\ell \in \lab(s)$, and
  $\lab(s) \in W(m)$ by \cref{def:ee-signature}.
\item An initialization belongs to the allocation of its object, which comes
  before it in the same block (\cref{sec:ee-fresh}). So its object was
  allocated during the invocation, and did not exist in $\sigma$.
\item An allocation writes $\alloc \in \Rest$, and $\Rest \in W(m)$.
\item A call of $m'$: $m' \in \callees(s)$ by \cref{as:ee-facts}, and $m'$ is
  not opaque, as $m$ is not. $\ell$ exists when the call starts, as it existed
  in $\sigma$, and the call's invocation writes it. Its derivation is a part of
  the derivation of $m$'s invocation, so by the induction hypothesis, $\ell$ lies
  in a partition of $W(m') \subseteq W(m)$.
\item Any other node: by \cref{as:ee-facts}. \qedhere
\end{itemize}
\end{proof}

Two things are outside the theorem. The first is the objects an invocation
allocates itself: their fields are not in the footprint, which is what lets the
facts leave initializations out (\cref{sec:ee-fresh}). The second is memory that
is not in $\Loc$. Graal also names locations of the virtual machine, such as the
pointers of a thread's allocation buffer, a thread's local variables, and
memory outside the Java heap. The facts list the ones a method writes, for the
kill sets of \cref{sec:ee-kills}, but neither the signatures nor the token form
model them, and the facts do not list the ones a method reads. So the tokens
order the operations on Java's memory: fields, array elements and the
allocator. They do not order a write and a read of memory outside the heap made
by two different callees. In the image of Are We Fast Yet, 25 methods write
such memory. No rewrite of this thesis moves a call, so none of them depends
on it, but a scheduler on the token form would need a token for the virtual
machine's memory too.

%----------------------------------------------------------------------------
\section{What the tokens guarantee}\label{sec:ee-guarantee}

The first guarantee is that an instruction that does not take a partition's
token leaves the partition alone. Only fields of objects allocated in the
meantime may change, and no operation before them can have seen those.

\begin{lemma}[Frame]\label{lem:ee-frame}
Let $s_j$ and $s_k$, $j < k$, be steps of a run of an invocation of $m$, and
$P \in T_m$, such that no step strictly between them takes $P$, except moves
and guards. Then every location of $P$ that exists after $s_j$ has the same
value in $\sigma_k$ as in $\sigma_{j+1}$.
\end{lemma}

\begin{proof}
Let $\ell \in P$ exist in $\sigma_{j+1}$. We show that no step $s_i$,
$j < i < k$, writes $\ell$. Moves, guards, operations on values, copies and
loads write nothing. A store or an initialization that takes the token of
another partition $Q$ writes a location of $Q$, by \cref{prop:ee-sound-labels},
and $Q \cap P = \emptyset$. An initialization that takes no token writes a
location whose partition is not in $T_m$, by \cref{def:ee-token-form}, so not
$\ell$, as $P \in T_m$. An allocation that does not take $P$ writes
$\alloc \notin P$ and fields of an object that did not exist before. A call that
does not take $P$ calls methods $m'$ with $P \notin T_{m'}$, none of them
opaque. $\ell$ exists when the call starts, so by \cref{thm:ee-signatures} the
call writes $\ell$ only if $\ell$ lies in a partition of
$W(m') \subseteq T_{m'}$, which it does not.
\end{proof}

The second guarantee is that two instructions that share no token can be
swapped. Its proof needs a property of the semantics that separation logic
calls the frame property~\cite{yang2002local,reynolds2002separation}: an
instruction's behaviour depends only on the locations it reads.

\begin{lemma}[Locality]\label{lem:ee-locality}
Let $\cfg{\iota, \rho, \sigma} \Downarrow r$, and let $\hat\sigma$ agree with
$\sigma$ on every location that exists in $\sigma$ and that the step reads,
and on $\alloc$ if it allocates. Then $\cfg{\iota, \rho, \hat\sigma}
\Downarrow \hat r$, where $\hat r = \Trap$ if $r = \Trap$, and if
$r = \cfg{\rho', \sigma'}$, then $\hat r = \cfg{\rho', \hat\sigma'}$ for the
heap $\hat\sigma'$ that agrees with $\sigma'$ on every location the step writes
or allocates, and with $\hat\sigma$ everywhere else.
\end{lemma}

\begin{proof}
By induction on the derivation of $r$. Each rule reads only the locations it
names, which the step reads; with the same values for them, every premise has
the same outcome, and the rule writes the same locations with the same values.
\end{proof}

We also need that a reference never points to an object that does not exist:
in every heap of a run, the fields of every existing object hold $\nul$, numbers
or existing objects, and so do the environments. This holds initially and every
rule preserves it.

\begin{theorem}[Independence]\label{thm:ee-independence}
Let $\iota_1\,\iota_2$ be adjacent instructions of a block of $m$, such that
$\iota_2$ uses no result of $\iota_1$, and the two take no common token. Then
from every environment and heap a run reaches before them, $\iota_1\,\iota_2$
and $\iota_2\,\iota_1$ have the same normal outcomes: one of them gives
$\cfg{\rho', \sigma'}$ if and only if the other does.
\end{theorem}

\begin{proof}
Let $\iota_1\,\iota_2$ go from $\cfg{\rho, \sigma}$ through $\sigma_1$ to
$\cfg{\rho', \sigma'}$, and let $T_1$ and $T_2$ be the tokens the two take.
First, consider an initialization that takes no token: its location lies in a
partition no other instruction of $m$ or of its callees touches, since a load,
store or call that touched it would put its partition into $T_m$
(\cref{def:ee-signature}), and the initialization would then take the token. So
it commutes with every instruction but the allocation of its object, which
gives $\iota_2$ a value it uses. Otherwise:
\begin{enumerate}
\item \emph{The footprints are disjoint.} By \cref{prop:ee-sound-labels} for
  accesses and \cref{thm:ee-signatures} for calls, every location that exists
  before a step and that it reads or writes lies in a partition the step takes.
  A location that exists in $\sigma$, and that one of them writes and the other
  reads or writes, would lie in a partition of $T_1 \cap T_2 = \emptyset$.
\item \emph{At most one of them allocates}, since both would take $\Rest$.
\item \emph{$\iota_2$ never reaches an object $\iota_1$ allocated.} Its
  arguments do not come from $\iota_1$, so they are $\nul$, numbers or objects
  that exist in $\sigma$. Every reference it reads later comes from a field of
  an object that exists in $\sigma$, which $\iota_1$ did not write by (1), so the
  field holds what it held in $\sigma$, an object that existed then; or from a
  field of an object $\iota_2$ allocated itself.
\item By (1)--(3), $\sigma$ and $\sigma_1$ agree on everything $\iota_2$ reads,
  and on $\alloc$ if $\iota_2$ allocates. By \cref{lem:ee-locality}, $\iota_2$
  from $\cfg{\rho, \sigma}$ gives the same values and the same writes, and
  reaches a heap $\sigma_2$.
\item In the same way, $\sigma$ and $\sigma_2$ agree on everything $\iota_1$
  reads from $\sigma$, so $\iota_1$ from $\sigma_2$ gives the same values and
  writes as from $\sigma$.
\end{enumerate}
Both orders bind the same values and make the same writes, to disjoint
locations, and allocate the same objects, so they end in the same environment
and heap. The converse is symmetric.
\end{proof}

\begin{corollary}[Any order]\label{cor:ee-any-order}
Any two orders of the instructions of a block that respect its data flow, of
values and of tokens alike, have the same normal outcomes.
\end{corollary}

\begin{proof}
The data flow orders the instructions partially, and two linear extensions of
a partial order differ by a sequence of swaps of adjacent incomparable
elements. Each swap preserves the normal outcomes by
\cref{thm:ee-independence}.
\end{proof}

\begin{remark}[Abnormal ends]\label{rem:ee-abnormal}
The theorem says nothing about traps and nontermination, and cannot. A guard
takes every token of its method, but a call of a method that touches no memory
takes none, so the two share no token. If the guard traps and the callee never
returns, swapping them turns a trap into a program that hangs. No rewrite of
this thesis moves an instruction, so none depends on it. A scheduler on the
token form would need termination, like memory, to be an effect with a
token.
\end{remark}

%----------------------------------------------------------------------------
\section{What the tokens license}\label{sec:ee-license}

Two rewrites follow from the guarantees: replacing a load by a value already
known, and removing a store nothing reads. For each, we first say what must hold
in every run, then what the token form shows.

\subsection{Forwarding a load}

\begin{definition}[Forwardable load]\label{def:ee-forwardable}
A load $\iota$, $x, h' \coloneqq \getf{f}{P}(h, y)$, of a method $m$ is
\emph{forwardable} from a variable $v$ bound before it in its block if in every
run of an invocation of $m$, every step $s_k$ that executes $\iota$ has an
earlier step $s_j$ such that
\begin{enumerate}
\item $s_j$ is a load, a store or an initialization of field $f$ of the object
  $o = \rho_k(y)$, and $\rho_k(v)$ is the value it read or wrote, and
\item no step strictly between $s_j$ and $s_k$ takes $P$, except moves and
  guards.
\end{enumerate}
\end{definition}

\begin{theorem}[Forwarding]\label{thm:ee-forwarding}
Replacing a forwardable load $\iota$ by $x \coloneqq \move\ v;\ h' \coloneqq
\move\ h$ preserves the outcome of the program, in the semantics of
\cref{fig:ee-semantics} and in Java's (\cref{sec:ee-java}).
\end{theorem}

\begin{proof}
Let $p'$ be the program with $\iota$ replaced. We show by induction on
derivations, in both directions, that every invocation has the same outcome in
$p$ and in $p'$. Invocations of other methods than $m$ differ only in the
invocations they make, so consider one of $m$, and its run in $p$. Up to a step
$s_k$ that executes $\iota$, $p'$ executes the same steps, with the same
environments and heaps. At $s_k$, $p$ reads $\sigma_k(o, f)$ with
$o = \rho_k(y)$. $o \neq \nul$, since $s_j$ accessed it, so $\iota$ does not
trap. $s_j$ left $\rho_k(v)$ in $(o, f)$, which exists after $s_j$ and lies in
$P$, since $\iota$ accesses it (\cref{prop:ee-sound-labels}). By
\cref{lem:ee-frame}, $\sigma_k(o, f) = \sigma_{j+1}(o, f) = \rho_k(v)$. So $x$
gets the same value in $p$ and $p'$, and neither changes the heap; from there
on the runs are the same again.
\end{proof}

The proof never mentions what happens when something traps: forwarding a load
does not depend on how exceptions are modelled.

\begin{lemma}[Straight tokens]\label{lem:ee-straight}
Let a load $\iota$ of $f$ through $y$ take a token of $P$ that comes, through
moves and guards only, from the token an access $\iota_0$ of $f$ through $y$ in
the same block gives. Then $\iota$ is forwardable from the value $\iota_0$
reads or writes.
\end{lemma}

\begin{proof}
By \cref{lem:ee-chains}, in every execution of the block, the steps between
$\iota_0$ and $\iota$ that take $P$ are the moves and guards on the token's
way. $y$ is bound once, so both access the same object.
\end{proof}

Between blocks, the rules of \cref{chap:analysis} first pass the value to the
load's block as a parameter, and then apply the lemma within the block
(\cref{sec:ee-shape}).

\begin{example}[The running example, forwarded]\label{ex:ee-running-forwarded}
Under object-field, the rules join the three blocks of
\cref{fig:ee-running-tokens} into one and turn the null check into a guard on
each token. The load's token then comes straight from the store, and both go
through \texttt{counter}, so by \cref{lem:ee-straight} the load is forwardable
from the stored 5. In the names of \cref{fig:ee-running-tokens}:
\begin{lstlisting}[language=ct]
    heap set_0 h_c_g p0 v11 → h_c1                    ; counter.value = 5
    Aliasing_countInto run → k13
    f_rrhh_hh call k13 p1 p2 h_o_g h_n_g → h_o1 h_n1  ; countInto(other, list)
    int move v11 → v24                                ; forwarded Aliasing$Counter.value
    heap move h_c1 → h_c2
\end{lstlisting}
Under field and with one partition, the load's token comes from the call, and
the lemma does not apply.
\end{example}

\subsection{Removing a store}

\begin{definition}[Overwritten store]\label{def:ee-overwritten}
A store $\iota$, $h' \coloneqq \setf{f}{P}(h, y, z)$, of a method $m$ is
\emph{overwritten} if in every run of an invocation of $m$, every step $s_j$
that executes $\iota$ has $o = \rho_j(y) \neq \nul$ and is followed by one of:
\begin{enumerate}
\item a step $s_k$ that stores to $(o, f)$, such that no step strictly between
  $s_j$ and $s_k$ takes $P$, except moves and guards;
\item the end of the run with $\Trap$, or no end at all, such that no step after
  $s_j$ takes $P$, except moves and guards.
\end{enumerate}
\end{definition}

In particular, the run does not return between a store that is overwritten and
the store that overwrites it.

\begin{theorem}[Dead stores]\label{thm:ee-dead}
Removing an overwritten store, that is, replacing it by $h' \coloneqq \move\ h$,
preserves the outcome of the program in the semantics of
\cref{fig:ee-semantics}.
\end{theorem}

\begin{proof}
Let $p'$ be the program without $\iota$, and consider a run of $m$ in $p$ that
executes $\iota$ at step $s_j$. Up to $s_j$, $p'$ executes the same steps. From
$s_{j+1}$ on, until the overwrite in case 1, the two runs execute the same
steps, with heaps that differ at $(o, f)$ alone. For that, it is enough that no
step between $s_j$ and $s_k$ reads or writes $(o, f)$, since by
\cref{lem:ee-locality} every such step then does the same in both. A step that
reads or writes $(o, f)$ takes $P$: a load or a store of it by
\cref{prop:ee-sound-labels}; a call by \cref{thm:ee-signatures}, since
$(o, f)$ exists when the call starts. An initialization that takes no token
belongs to the allocation of its object, and $o$ was allocated before $s_j$.
In case 1, $s_k$ then stores the same value to $(o, f)$ in both runs, after
which their heaps are equal, and the rest of the runs are the same. In case 2,
both runs end with $\Trap$, or neither ends. The converse is symmetric.
\end{proof}

The second case is where the model of exceptions comes in. In the semantics of
\cref{fig:ee-semantics}, a path that throws ends the program with $\Trap$,
which carries no heap, so nothing reads the store on it: this is the \emph{trap
model}. It is how the translation sees a path that throws, since it turns every
block from which every path throws into a trap (\cref{chap:conversion}).

\subsection{In Java}\label{sec:ee-java}

In Java an exception does not end the program. A handler up the stack may catch
it and read memory, and so may the code that prints an uncaught exception's
message. The \emph{Java model} changes the semantics in two ways:
\begin{itemize}
\item a guard that holds, an access through $\nul$ and $\ktrap$ give an exception
  $\Exc(\sigma)$, which carries the heap, and a call passes on the exception
  its callee gives;
\item any call may throw \texttt{StackOverflowError}, and any allocation
  \texttt{OutOfMemoryError}, so a call or an allocation may give
  $\Exc(\sigma)$ at any time, before it does anything.
\end{itemize}
The semantics is no longer deterministic, and the outcome of an invocation is
the set of the outcomes $\Ret(\bar v, \sigma)$ and $\Exc(\sigma)$ it may give.

\begin{definition}[Overwritten in Java]\label{def:ee-overwritten-java}
A store is \emph{overwritten in Java} if it is overwritten by the first case of
\cref{def:ee-overwritten} alone, and no step strictly between $s_j$ and $s_k$
may throw: none is a call, an allocation, a guard or an access through $\nul$.
\end{definition}

\begin{theorem}[Dead stores in Java]\label{thm:ee-dead-java}
Removing a store that is overwritten in Java preserves the outcomes of the
program in the Java model.
\end{theorem}

\begin{proof}
As for \cref{thm:ee-dead}, the two runs differ at $(o, f)$ alone between $s_j$
and $s_k$. No step there may throw, so no exception carries a heap in which the
two differ: every exception is thrown before $s_j$ or after $s_k$, where the
heaps agree. A run that enters the window either reaches $s_k$ in both programs
or runs forever in both.
\end{proof}

\begin{example}[Legal in Java, or only in the trap model]\label{ex:ee-java}
In \texttt{deadStorePastOtherObject} of \texttt{Aliasing},
\begin{lstlisting}[language=Java]
int seen = other.other;
counter.value = 1;          // dead
seen += other.value;
counter.value = 2;
\end{lstlisting}
under object-field the load of \texttt{other.value} takes the token of
\texttt{other}'s counters, so the first store's token goes straight to the
second store. Apart from arithmetic, only the load comes between the two
stores, and its receiver was checked for $\nul$ before the first store, by the
load of \texttt{other.other}. So nothing between them may throw, and the store
is dead in Java. In \texttt{deadStoreReadInArm} of \texttt{Examples}, once the load in a
branch arm takes the value of the first store, the second store overwrites it,
and the call \texttt{length(list)} comes between them. Under field the call
takes only the token of \texttt{Link.next}, so the store is overwritten, but the
call may throw \texttt{StackOverflowError}, and a handler could see the first
value: the store is dead in the trap model only.
\end{example}

The implementation removes a store in Java only if nothing between it and its
overwrite may throw, on any token: a call, an allocation, a division, a guard or
a trap (\cref{chap:analysis}). Those are the stores it gives Native Image to
remove. The trap model's are only counted.

\subsection{Finer partitions license more}\label{sec:ee-monotone}

The rules check the conditions of this section on the token form, by following
tokens.

\begin{definition}[Static conditions]\label{def:ee-static}
In the token form of a method, a load is \emph{straight} if it satisfies the
condition of \cref{lem:ee-straight}. A store of $f$ through $y$ that takes $P$
is \emph{statically overwritten} if every \emph{$P$-path} from it ends at a store
of $f$ through $y$, or at a $\ktrap$. A $P$-path follows the control flow from
the store: within a block to the next instruction that takes $P$ and is not a
move or a guard, or to the terminator; through a jump or either arm of a branch
to the start of the next block, following $y$ into the parameter it is passed
to. It ends there, or at a $\kret$. A path that comes back to a block it went
through, with the token and $y$ in the same parameters, is cut. A store is
\emph{statically overwritten in Java} if, moreover, no $P$-path ends at a
$\ktrap$ or goes through a guard, and none of the instructions the paths go
past may throw.
\end{definition}

A run that executes a statically overwritten store follows one of the
$P$-paths of the store. By \cref{lem:ee-chains}, the steps in which it takes
$P$ are the instructions on that path. A run that goes around a cut loop forever takes $P$
in no step after the store. So a statically overwritten store whose receiver
is never $\nul$ is overwritten, and in Java too.

\begin{theorem}[Monotonicity]\label{thm:ee-monotone}
Let $\Pi' \sqsubseteq \Pi$ be two of the four partitionings. Every load that is
straight under $\Pi$ is straight under $\Pi'$, and every store statically
overwritten under $\Pi$, in the trap model or in Java, is statically
overwritten under $\Pi'$.
\end{theorem}

\begin{proof}
First, every instruction that takes a token of $P' \in \Pi'$ in the token form
under $\Pi'$ takes the token of $\up{P'}$ under $\Pi$:
\begin{itemize}
\item a load or a store by \cref{prop:ee-order}; an initialization too, since
  if its label $P'$ is in $T'_m$, then $\up{P'}$, its label under $\Pi$, is in
  $T_m$, as for a call below;
\item an allocation, since $\Rest$ is the same partition in both;
\item a call, a guard or another node, which take the tokens of a method $n$,
  its callee or its own: if $P' \in T'_n$, then either $P'$ is the label under
  $\Pi'$ of an access in a method $n$ reaches, and $\up{P'}$ is that access's
  label under $\Pi$, which is in $T_n$; or $n$ is opaque, under $\Pi$ as well,
  since opacity does not depend on the partitioning.
\end{itemize}
Now let a load $\iota$ of $f$ through $y$ be straight under $\Pi$, from
$\iota_0$, with label $P'$ under $\Pi'$ and $P = \up{P'}$ under $\Pi$.
$\iota_0$ accesses $f$ through $y$ too, so its label under $\Pi'$ is also $P'$.
Under $\Pi$, the instructions between them that take $P$ are moves and guards.
An instruction between them that takes $P'$ under $\Pi'$ takes $P$ under $\Pi$,
so it is a move or a guard, and $\iota$ is straight under $\Pi'$.
For a store, the $P'$-paths and the $P$-paths follow the same control flow,
and on each, the instructions that take $P'$ are among those that take $P$. The
first one that is not a move or a guard on a $P$-path is a store of $f$ through
$y$, which takes $P'$ too, so it is the first on the $P'$-path as well; or the
path reaches a $\ktrap$ or is cut, and so does the $P'$-path. The conditions
added in Java do not depend on the partitioning: guards take every token of
their method under both.
\end{proof}

The theorem is about the conditions. The rules apply them in some order,
interleaved with changes of shape that copy code (\cref{sec:ee-shape}), and the
copies count against bounds: peeling a loop copies at most 40 instructions of
its header and body. A finer partitioning has more tokens, so a null check
becomes more guards, and a loop a coarser partitioning peels may be too long to
peel under a finer one. So the implementation is not monotone by construction,
and we checked that it is on the programs of this thesis. \Cref{tab:ee-monotone}
shows, for each program, what the rules leave under the four partitionings,
all from the export with allocation sites, so that the program is the same
under each. In no method does a finer partitioning leave more stores or more
loads than a coarser one, in the whole method or in its loops, and every store
that is dead in Java under a coarser partitioning is dead under a finer one.
A check of the test suite does the same on the three programs it uses.

\begin{table}[t]
\centering
\begin{tabular}{@{}lrrrrr@{}}
\toprule
Program & Methods & none & field & object & object-field \\
\midrule
\texttt{Aliasing} & 14 & 18/15 & 18/15 & 17/12 & 17/12 \\
\texttt{Examples} & 41 & 51/37 & 44/29 & 45/29 & 44/29 \\
\texttt{Signatures} & 18 & 10/10 & 10/6 & 10/6 & 10/6 \\
Are We Fast Yet & 370 & 223/345 & 223/308 & 223/321 & 223/308 \\
the native-image driver & 427 & 209/675 & 209/656 & 209/665 & 209/656 \\
\bottomrule
\end{tabular}
\caption{The field stores and loads the rules leave in the translated methods,
under each partitioning, on the export with allocation sites. No method keeps
more under a finer partitioning than under a coarser one.}
\label{tab:ee-monotone}
\end{table}

\subsection{From licences to rules}\label{sec:ee-shape}

The rules of \cref{chap:analysis} do more than the two lemmas: they forward
into the arms of a branch, into and out of loops, and across merges. They do it
by changing the shape of the program first, so that the value reaches the
load's block as a parameter, and then applying \cref{lem:ee-straight} within the
block.

\begin{lemma}[Shapes]\label{lem:ee-shapes}
Each of the following changes preserves the outcome of the program, and for
every partition $P$, the sequence of steps that take $P$ other than moves and
guards, with their values:
\begin{enumerate}
\item a branch to a block that traps on every path becomes a guard on each token
  and a jump to the other block;
\item a block that only one jump reaches joins the block that jumps;
\item a parameter of a block for which every jump and branch passes the same
  value as for another parameter is replaced by that one;
\item a block takes a new parameter, which every jump and branch to it passes;
\item the first iteration of a loop is peeled: copies of its header and body run
  it, and jump into the loop;
\item a merge is split: the blocks that jump to it are divided in two, one of
  them jumps to a copy of the merge, and both copies go on in a new block with
  the code after a given point.
\end{enumerate}
\end{lemma}

\begin{proof}[Proof sketch]
In (1), the block that traps ends the run with $\Trap$ after no step that
touches memory, and the guard does too, at the same point; otherwise both go on
in the other block. Changes (2)--(6) rename and copy blocks: every run of the
new program executes the same instructions, in the same order and with the same
values, as a run of the old one, and the other way round.
\end{proof}

A rule then forwards a load by passing, as the new parameter of (4), the value
of the access each jump's token comes from. In every run the parameter then
holds the value of the access before the load, and the load is forwardable from
it. For a load in a loop, this is the value the entry knows, or the value the
body left the iteration before, a loop-carried $\varphi$.

%----------------------------------------------------------------------------
\section{Fresh objects}\label{sec:ee-fresh}

Graal gives the objects that escape analysis keeps virtual their first values
without stores. When such an object escapes, a single node, the commit of an
allocation, allocates it and gives each of its fields its first value. Graal's
model of memory has a location of its own for such writes,
\texttt{LocationIdentity.init()}, which no read ever names: a field of a new
object can only be read through a reference to it, and the reference only
exists once the commit is done.

The translation writes a commit as the allocations followed by an
initialization of each field, and the facts leave initializations out of the
signatures. So an initialization takes a token only when its method takes the
token of the field's partition for another reason, and otherwise takes none.

\begin{proposition}[Fresh objects]\label{prop:ee-fresh}
The results of this chapter hold for the token form with initializations that
take no token.
\end{proposition}

\begin{proof}
Each proof treats them as a case of its own: \cref{thm:ee-signatures}, because
the object did not exist when the invocation started; \cref{lem:ee-frame},
because the partition is not one of the method's; \cref{thm:ee-independence},
because nothing else in the method or its callees touches the partition; and
\cref{thm:ee-dead}, because the stored object was allocated before the store.
\end{proof}

This is also why \cref{lem:ee-frame} is about the locations that exist after
the earlier step. A callee may allocate an object, initialize a field of it in
a partition of the caller, and give it back; that field then changes between two
steps of the caller that take the partition, without the call taking it. But no
earlier operation can have seen the field, so no rewrite based on one can be
about it.

\Cref{thm:ee-signatures} constrains only locations that exist when an
invocation starts. So an access that, in every run, touches only fields of
objects allocated during the current invocation of its method could be left out
of the signature altogether: no caller can tell it happened. The method itself
still has to order such accesses among themselves, with a token of its own that
it creates where it allocates the objects and consumes before it returns, as
Haskell's state threads do~\cite{launchbury1994state}. An initialization needs
no such token, because when it takes none, nothing else in its method touches
its partition. Leaving fresh objects out of a signature is the \emph{prestate}
of S\u{a}lcianu and Rinard's purity analysis~\cite{salcianu2005purity}, and the
effect masking of Lucassen and Gifford's effect
systems~\cite{lucassen1988effects}. Telling which accesses touch only fresh
objects takes an analysis of its own: an allocation site is not enough, since it
allocates in every invocation, while S\u{a}lcianu and Rinard's analysis tells
apart the objects allocated during the current one. The implementation uses
freshness only where Graal already knows it, for commits.

%----------------------------------------------------------------------------
\section{Concurrency}\label{sec:ee-concurrency}

The semantics of \cref{fig:ee-semantics} is sequential: a run is one thread.
Java threads share the heap, and another thread may write a location between
two steps of a run. The Java memory model~\cite{manson2005jmm} promises
sequential consistency to programs free of data races~\cite{adve1990weak}: if no
sequentially consistent execution of a program has two accesses to one location
from two threads, at least one of them a write, that synchronization does not
order, then every execution behaves as an interleaving of the steps of its
threads. For such programs, removing a load that an earlier access to the same
location in the same thread makes redundant, and a store that a later store
overwrites, is safe when no synchronization comes between
them~\cite{sevcik2011safe}. For programs with data races, the guarantees of the
Java memory model are weaker, and which transformations it allows is
subtle~\cite{sevcik2008validity}. We make no claims about such programs.

In the facts, every synchronization (a lock, a volatile access, an atomic
update, a fence) makes its method opaque, so a call that synchronizes takes
every token, and nothing is forwarded or removed across it. The translation
refuses methods that make a volatile access themselves (\cref{chap:conversion}).

That is more than necessary. Synchronization makes the writes of other threads
visible, and other threads can only write locations they can reach.

\begin{proposition}[Thread-local partitions]\label{prop:ee-thread-local}
Let the program be free of data races, and let $P$ be a partition whose
locations no two threads access in any execution. Then \cref{lem:ee-frame}, and
with it the results that rely on it, hold for $P$ even if a step that only
synchronizes does not take $P$.
\end{proposition}

\begin{proof}[Proof sketch]
Every execution is an interleaving of the steps of its threads. Steps of other
threads access no location of $P$, and a synchronization touches its own lock
or volatile field, not a location of $P$. So between two steps of a run, only
steps of the same run write the locations of $P$, and the proof of
\cref{lem:ee-frame} applies to them.
\end{proof}

Which partitions are thread-local is what a thread-escape analysis
tells~\cite{choi1999escape,salcianu2001escape}. That the program starts no thread
is not enough: Native Image runs threads of its own, which handle references the
program registers, such as cleaners. In Are We Fast Yet, after the kill sets
were made precise, a lock is the first effect 333 of the 455 calls that still
kill everything reach, 294 of them through the lock
\texttt{Throwable.fillInStackTrace} takes on the exception it is building,
which no other thread can see yet. Run-time class initialization, which runs
code under a lock, accounts for 108 more (\cref{chap:evaluation}). This
condition is not implemented.

%----------------------------------------------------------------------------
\section{What a kill set guarantees}\label{sec:ee-kills}

Signatures need not go through tokens. Graal's own model of memory orders a read
after the last node before it that kills the read's location, and takes a call
to kill every location. The round trip of \cref{chap:conversion} gives each call
a \emph{kill set} instead: the locations its callees may write, made from the
same facts.

Graal's locations are the fields, one location per field or, under
object-field, per partition of a field; the array element kinds; $\mathsf{init}$;
the locations of the virtual machine; and $\mathsf{any}$, which overlaps all of
them. Every node that writes memory says which locations it kills.

\begin{assumption}[Graal's model]\label{as:ee-graal}
A node writes only locations it kills. The writes that initialize the objects a
node allocates are kills of $\mathsf{init}$.
\end{assumption}

Graal's own phases rely on this for every node, so we rely on it too.

\begin{definition}[Kill set]\label{def:ee-kill-set}
The kill set $K(m)$ of a method is $\mathsf{any}$ if $m$ is opaque. Otherwise it
is the least set that contains the locations of $m$'s own writes, namely the
field, or its partition, of a store; the element kind of an array store; the
locations another node kills; and $\mathsf{init}$ if $m$ allocates, together with
$K(m')$ for every callee $m'$ of every call in $m$. A call of $m$ kills
$K(m) \cup V$, where $V$ holds $\mathsf{init}$ and the locations of the virtual
machine that may change at any safepoint: the pointers of the thread's
allocation buffer, the salt of identity hash codes, the safepoint counter, and
the state of memory sessions.
\end{definition}

\begin{theorem}[Kill sets]\label{thm:ee-kills}
Assume \cref{as:ee-points-to,as:ee-facts,as:ee-graal}, and that the image never
executes a call its analysis found unreachable. Then every location written
during an invocation of $m$ is covered by a location of $K(m) \cup V$, where a
field of an object allocated during the invocation and given its first value by
a commit counts as covered by $\mathsf{init}$.
\end{theorem}

\begin{proof}[Proof sketch]
By induction on the derivation, as for \cref{thm:ee-signatures}, with Graal's
locations in place of partitions: a store writes its field's location, or its
partition's, which is in $K(m)$; another node writes what it kills, by
\cref{as:ee-graal}; a call writes what its callee does, which the analysis
lists, unless the call is unreachable, and the image then never makes it. What
changes at a safepoint is written by the virtual machine, not by the callee:
any call may reach a safepoint, and while it waits there, another thread may
collect garbage and reset the allocation buffer. Native Image declares the same
locations $V$ at its own safepoints.
\end{proof}

A kill set covers writes only. A read in the callee changes nothing that a
read in the caller could see, so the caller may still take a value from before
the call. What a read in the callee does need is that a store before the call
is neither removed nor moved past it. Graal keeps stores and calls in their
order, and the only rule with which we saw it remove a dead store requires the
two stores to be adjacent (\cref{chap:conversion}). So an effect that only reads
makes its method take every token, but leaves its kill set alone. The kill set
is to a call what $W(m)$ is to its tokens: a call takes the tokens of
$R(m) \cup W(m)$, but kills only the locations of $W(m)$.

%----------------------------------------------------------------------------
\section{Summary}\label{sec:ee-summary}

\begin{table}[t]
\centering
\small
\begin{tabular}{@{}>{\raggedright\arraybackslash}p{0.28\textwidth}>{\raggedright\arraybackslash}p{0.3\textwidth}>{\raggedright\arraybackslash}p{0.34\textwidth}@{}}
\toprule
Result & Assumes & Relied on by \\
\midrule
Chains (\ref{lem:ee-chains}) & single use, well typed & everything below \\
Signatures (\ref{thm:ee-signatures}) & sound points-to and facts & the frame, independence, dead stores \\
Frame (\ref{lem:ee-frame}) & signatures & forwarding, dead stores \\
Independence (\ref{thm:ee-independence}) & frame, locality & that tokens order what must be ordered \\
Forwarding (\ref{thm:ee-forwarding}) & frame & the rules' forwarding \\
Dead stores (\ref{thm:ee-dead}) & frame, the trap model & the rules' dead stores \\
Dead stores in Java (\ref{thm:ee-dead-java}) & nothing between may throw & the stores Native Image removes \\
Monotonicity (\ref{thm:ee-monotone}) & refinement & choosing a partitioning \\
Fresh objects (\ref{prop:ee-fresh}) & Graal's commits & signatures without initializations \\
Thread-local partitions (\ref{prop:ee-thread-local}) & no data races, a thread-escape analysis & not implemented \\
Kill sets (\ref{thm:ee-kills}) & Graal's model, the closed world & the round trip \\
\bottomrule
\end{tabular}
\caption{The results of this chapter.}
\label{tab:ee-summary}
\end{table}

\Cref{tab:ee-summary} lists the results. Stating them precisely, and holding
them against the implementation, turned up four things the earlier chapters
take for granted:
\begin{itemize}
\item \emph{The frame holds only for locations that already exist.} Between two
  operations on a partition, a callee may initialize a new object's field in it
  without taking its token. Nothing can have seen the field before, so no
  rewrite is affected, but the tokens do not carry the heap: an interpreter that
  gave each token its partition's contents would be wrong.
\item \emph{Tokens model Java's memory, not the virtual machine's.} The facts
  list the virtual machine's locations a method writes only for the kill sets,
  and not the ones it reads. No rewrite moves a call, so none depends on it.
\item \emph{Independence holds for normal outcomes only.} Tokens do not order a
  guard that may trap before a call that may not return. No rewrite moves an
  instruction, so none depends on it.
\item \emph{Monotonicity holds of the conditions, not of the rules by
  construction}, since peeling is bounded by a number of instructions that
  grows with the number of tokens. It holds on every program of this thesis,
  and a check keeps it so.
\end{itemize}
